\documentclass[10pt]{article}

% --- Page geometry ---
\usepackage[margin=1in]{geometry}

% --- Encoding and fonts ---
\usepackage[T1]{fontenc}
\usepackage[utf8]{inputenc}
\usepackage{lmodern}

% --- Math ---
\usepackage{amsmath,amssymb,amsthm}

% --- Figures and tables ---
\usepackage{graphicx}
% \figinclude{path}: include figure if file exists, else show placeholder box
\newcommand{\figinclude}[1]{%
  \IfFileExists{#1}%
    {\includegraphics[width=\linewidth]{#1}}%
    {\begin{center}\fbox{\parbox{0.7\linewidth}{\centering
      \textit{[Figure pending --- run analysis/figures.py]}
    }}\end{center}}%
}
\usepackage{booktabs}
\usepackage{multirow}
\usepackage{caption}
\usepackage{subcaption}
\usepackage{float}

% --- Colors and hyperlinks ---
\usepackage[dvipsnames]{xcolor}
\usepackage[colorlinks=true,linkcolor=MidnightBlue,citecolor=MidnightBlue,urlcolor=MidnightBlue]{hyperref}

% --- Bibliography ---
\usepackage[numbers,sort&compress]{natbib}

% --- Misc ---
\usepackage{enumitem}
\usepackage{algorithm}
\usepackage{algpseudocode}
\usepackage{xspace}

% --- Theorem environments ---
\newtheorem{proposition}{Proposition}
\newtheorem{theorem}{Theorem}

% --- Custom commands ---
\newcommand{\cate}{\tau(X)}
\newcommand{\EE}{\mathbb{E}}
\newcommand{\stpu}{\texttt{STPU}\xspace}
\newcommand{\mstpu}{\texttt{MSTPU}\xspace}
\newcommand{\causalforest}{\texttt{CausalForest}\xspace}
\newcommand{\rlearner}{\texttt{R-Learner}\xspace}
\newcommand{\tlearner}{\texttt{T-Learner}\xspace}
\newcommand{\ghat}{\hat{g}(X)}
\newcommand{\tauhat}{\hat{\tau}(X)}

% --- Title ---
\title{%
  \textbf{Can We Fix Wasted Ad Spend by Re-ranking Users?}\\[4pt]
  \large Sure-Thing Penalized Uplift: A Study of What Works and What Doesn't
}

\author{%
  Nathan Clark\\
  Independent Researcher\\
  \texttt{nathan.clark.dev@gmail.com}
}

\date{}

% ============================================================
\begin{document}
\maketitle

% ============================================================
\begin{abstract}
When you run an ad campaign, a big chunk of your budget goes to users who
would have bought anyway---they don't need the ad at all.
Our earlier work \citep{clark2025benchmarking} showed that all standard models
waste 86--92\% of their targeting budget this way, and that better model
accuracy doesn't fix it.

This paper tests a natural fix: lower the score of users who already have a
high chance of buying without the ad.  We call this Sure-Thing Penalized Uplift
(STPU).  We tested two versions:
\textbf{Additive STPU}---subtract a penalty from the score---and
\textbf{Multiplicative STPU} (MSTPU)---multiply the score by how much ``room''
the user has to be influenced.

The short version: neither works.

Additive STPU always makes things worse.  The reason is a simple numbers
problem: the penalty we subtract is 2--5$\times$ larger in scale than the
scores themselves, so the penalty takes over and we end up targeting the wrong
people entirely.  Multiplicative STPU fixes that scale problem but still shows
no consistent improvement across 567 experiments---it's basically a coin flip.

The deeper reason neither approach works is that at a 5\% conversion rate,
fewer than 7\% of users actually benefit from an ad on any given day.
Even a perfect model could only reduce budget waste to 78\%; the best real
models sit at 86--92\%.  That remaining gap can't be closed by re-ranking,
because there just isn't enough signal in binary (yes/no) conversion data to
reliably identify those rare responsive users.

Fixing budget waste requires richer data---like revenue per user or multiple
touchpoints---not smarter sorting of noisy binary signals.
\end{abstract}

% ============================================================
\section{The Problem We're Trying to Solve}
\label{sec:intro}

Imagine you're deciding which 30\% of users to show an ad to.
The goal is to pick the users who will actually buy \emph{because} of the ad,
not users who would buy regardless.

The people who would buy regardless are called \textbf{sure-things}.
Showing them an ad wastes money with zero upside.
A user who only buys when shown an ad is a \textbf{persuadable}---the ideal target.
And a user who is \emph{less} likely to buy after seeing an ad is a
\textbf{sleeping dog}---actively harmful to target.

Uplift models try to estimate the \emph{incremental effect} of showing an ad:
\begin{equation}
  \tau(X) = \EE\bigl[\text{buy with ad} - \text{buy without ad} \;\big|\; \text{user features } X\bigr]
\end{equation}
This is called the Conditional Average Treatment Effect (CATE).
A persuadable has $\tau > 0$; a sure-thing has $\tau \approx 0$;
a sleeping dog has $\tau < 0$.

\paragraph{The core failure.}
Our earlier paper \citep{clark2025benchmarking} showed that even the best
uplift models---including the state-of-the-art Causal Forest---waste 86--92\%
of their top-30\% targeting budget.  Worse, making the model more accurate
(lower PEHE error) doesn't fix this.  Sure-things have high predicted
conversion rates, and models keep mistaking ``likely to buy'' for
``incrementally valuable.''

\paragraph{The proposed fix.}
Here's the intuition: if a user has a high chance of buying even without the ad
(call this $\hat{g}(X)$, estimated from people who \emph{weren't} shown ads),
they're probably a sure-thing.  So penalize their score:
\begin{equation}
  \text{score}(X) = \tauhat - \beta \cdot \ghat
\end{equation}
This is STPU. We test whether it works, why it fails, and what to do instead.

\paragraph{What this paper contributes.}
\begin{enumerate}[topsep=2pt,itemsep=1pt]
  \item We show additive STPU always makes things worse, and explain why
    (a scale mismatch---\S\ref{sec:failure}).
  \item We propose and test a scale-invariant alternative, Multiplicative STPU
    (\S\ref{sec:mstpu}).
  \item We show that neither version works across 567 experiments and explain
    the fundamental limit that prevents any re-ranking approach from working here
    (\S\ref{sec:results}).
\end{enumerate}

% ============================================================
\section{Setup and Metrics}
\label{sec:background}

\paragraph{How we measure waste.}
We simulate running an ad campaign and targeting the top 30\% of users by
predicted score.  Then we check: what fraction of those targeted users had
a true incremental effect near zero (i.e., they were non-incrementals)?
That's \textbf{wasted spend}.  Lower is better.

We also track \textbf{iROAS efficiency}: the average true lift of targeted
users, expressed as a fraction of what a perfect oracle would achieve.
Higher is better. An iROAS of 0.5 means we're getting half the value we
could with perfect information.

\paragraph{Baseline conversion probability.}
$\hat{g}(X)$ is our estimate of ``how likely is this user to buy without
an ad?''  We estimate it by training a classifier on just the control group
(users who were never shown an ad).  Sure-things have high $\hat{g}$;
persuadables have low $\hat{g}$.  This is the signal STPU tries to use.

\paragraph{What we tested.}
We ran experiments on three synthetic datasets (linear effects, nonlinear
effects, and mixed subgroups), three model families (Causal Forest, R-Learner,
T-Learner), at five confounding levels, using 20,000 users per run.
All told: 567 experiment configurations.

% ============================================================
\section{Why Additive STPU Fails: A Scale Problem}
\label{sec:failure}

\subsection{The idea behind it}

The additive STPU score is:
\begin{equation}
  s(X;\beta) = \underbrace{\tauhat}_{\text{estimated lift}} -
               \beta \cdot \underbrace{\ghat}_{\text{organic conversion rate}}
  \label{eq:additive}
\end{equation}
The logic: if the model overestimates lift for sure-things (because it
confuses ``they buy a lot'' with ``they respond to ads''), subtracting
their organic rate should cancel that out.

\subsection{The problem: the numbers are on completely different scales}

Here's the catch.  With a 5\% base conversion rate:
\begin{itemize}[topsep=2pt,itemsep=1pt]
  \item $\hat{g}(X)$ is a probability between 0 and 1.
    Across users, it typically has a standard deviation of 0.05--0.10.
  \item $\hat{\tau}(X)$ is the \emph{difference} between two probabilities.
    For rare events, this is much smaller---standard deviation
    around 0.02--0.04, even for the best models.
\end{itemize}

That means $\hat{g}$ is 2--5$\times$ ``louder'' than $\hat{\tau}$.
At $\beta = 0.5$, the penalty term $\beta \cdot \hat{g}$ has roughly
$4\times$ the variance of $\hat{\tau}$.  The penalty drowns out the
signal, and we end up ranking users by \emph{lowest organic conversion}
rather than highest lift.  That's basically backwards.

\paragraph{The data backs this up.}
Across every configuration we tested, additive STPU at any $\beta > 0$
made wasted spend \emph{worse}---never better.
Figure~\ref{fig:additive_failure} shows the effect of increasing $\beta$:
wasted spend climbs monotonically, no matter which DGP or model we use.

\begin{figure}[t]
  \centering
  \figinclude{figures/fig2_additive_failure.pdf}
  \caption{Wasted spend and iROAS efficiency as we increase the penalty $\beta$
    (Causal Forest, moderate confounding, 5\% conversion rate).
    As $\beta$ increases, wasted spend gets steadily worse across all datasets.
    iROAS can improve slightly because the penalty pushes out sleeping dogs
    (users who are actually hurt by ads), but this comes at the cost of also
    pushing out persuadables.}
  \label{fig:additive_failure}
\end{figure}

% ============================================================
\section{A Better Formulation: Multiplicative STPU}
\label{sec:mstpu}

\subsection{The idea}

Instead of subtracting $\hat{g}$, we \emph{multiply} $\hat{\tau}$ by
how much room there is to influence the user:
\begin{equation}
  s(X) = \tauhat \cdot \bigl(1 - \ghat\bigr)
  \label{eq:multiplicative}
\end{equation}
We call this Multiplicative STPU, or \mstpu.

Think of $(1 - \hat{g})$ as a ``persuadability weight.''  A user who
already has a 90\% chance of buying gets a weight of 0.10---their score
gets cut to 10\% of what it was.  A user with a 5\% organic conversion
rate gets a weight of 0.95---barely changed.

\subsection{Why this fixes the scale problem}

$(1-\hat{g})$ is always between 0 and 1, so it acts as a proportional
shrinkage, not an additive offset.  The sign and rough magnitude of the
score stay the same as $\hat{\tau}$.  No $\beta$ to tune.

Mathematically, for a sure-thing ($\tau \approx 0$, $\hat{g} \approx$ high):
$s = 0 \times \text{something} \approx 0$---correctly depressed.
For a persuadable ($\tau > 0$, $\hat{g} \approx$ low):
$s \approx \hat{\tau} \times 0.95 \approx \hat{\tau}$---barely changed.

\subsection{How to use it}

\begin{algorithm}[H]
\caption{Multiplicative STPU in practice}
\begin{algorithmic}[1]
  \Require Training data, a base uplift model
  \State Train your uplift model normally; get predicted lift $\hat{\tau}(X)$ for each user
  \State Train a separate classifier on the \emph{control group only} (no-ad users);
         get predicted organic conversion $\hat{g}(X)$
  \State Multiply: $\text{score} = \hat{\tau}(X) \times (1 - \hat{g}(X))$
  \State Target the top $k$\% by this new score
\end{algorithmic}
\end{algorithm}

% ============================================================
\section{Results: Neither Approach Consistently Helps}
\label{sec:results}

\subsection{The headline number}

Multiplicative STPU improves iROAS efficiency in 48\% of our 567 experiments.
That's statistically the same as a coin flip.

The typical improvement when it does help: $+0.001$ to $+0.003$ percentage
points in iROAS.  That's below the noise level across random seeds.
Figure~\ref{fig:main_result} shows the distribution of gains---mixed positive
and negative with no clear pattern.

\begin{figure}[t]
  \centering
  \figinclude{figures/fig1_main_result.pdf}
  \caption{The iROAS improvement from switching to \mstpu\ vs.\ the base model.
    Positive bars mean \mstpu\ helped; negative bars mean it hurt.
    The pattern is essentially random---no dataset or model shows consistent gains.}
  \label{fig:main_result}
\end{figure}

\subsection{Wasted spend never improves}

Table~\ref{tab:main_results} shows wasted spend side-by-side for base models
and \mstpu.  In every row, \mstpu\ leaves wasted spend flat or slightly worse.

Why?  At a 5\% conversion rate, only about 6.5\% of users have realized a
positive individual treatment effect (i.e., they actually would have bought
because of the ad on this particular day).  Even a perfect oracle, with full
knowledge of every user's true effect, still wastes 78\% of its targeting
budget---because 22\% of the users it could target have positive expected effects,
but only 6.5\% of all users showed realized $\tau = +1$ in the data.

Real models sit at 86--92\% waste.  The gap between oracle (78\%) and reality
(86--92\%) exists because of poor CATE estimation quality, not poor ranking.
Re-ranking what we already estimated can't close a gap that was created by
estimation error.

\begin{table}[t]
\centering
\caption{Benchmark results: base model vs.\ Multiplicative STPU ($\alpha=1.0$, $\rho=0.05$, mean over 5 seeds). $\downarrow$ better for wasted spend; $\uparrow$ better for iROAS.}
\label{tab:main_results}
\small
\begin{tabular}{lll rr rr rr}
\toprule
DGP & Model & Lrnr & \multicolumn{2}{c}{Wasted Spend $\downarrow$} & \multicolumn{2}{c}{iROAS Eff.\ $\uparrow$} & \multicolumn{2}{c}{Qini $\uparrow$} \\
\cmidrule(lr){4-5}\cmidrule(lr){6-7}\cmidrule(lr){8-9}
& & & Base & MSTPU & Base & MSTPU & Base & MSTPU \\
\midrule
\multicolumn{9}{l}{\textit{Linear CATE}} \\
  & CausalForest & RF & 0.921 & $0.922$ & 0.050 & $0.050$ & -0.099 & -0.098 \\
  & R-Learner & LR & 0.912 & $0.914$ & 0.095 & $\mathbf{0.095}$ & 0.596 & 0.607 \\
  & T-Learner & RF & 0.919 & $0.920$ & 0.088 & $0.082$ & 0.627 & 0.614 \\
\midrule
\multicolumn{9}{l}{\textit{Nonlinear CATE}} \\
  & CausalForest & RF & 0.929 & $0.932$ & 0.022 & $0.015$ & -0.535 & -0.545 \\
  & R-Learner & LR & 0.916 & $0.921$ & 0.031 & $0.026$ & 0.071 & 0.009 \\
  & T-Learner & RF & 0.923 & $0.924$ & -0.004 & $\mathbf{-0.002}$ & -0.082 & -0.088 \\
\midrule
\multicolumn{9}{l}{\textit{Heterogeneous CATE}} \\
  & CausalForest & RF & 0.903 & $0.903$ & 0.006 & $\mathbf{0.009}$ & 1.199 & 1.186 \\
  & R-Learner & LR & 0.913 & $0.915$ & 0.025 & $0.021$ & -3.053 & -2.980 \\
  & T-Learner & RF & 0.898 & $0.899$ & 0.020 & $0.017$ & 1.400 & 1.197 \\
\midrule
\bottomrule
\end{tabular}
\end{table}

\subsection{What \mstpu\ actually does to the targeting set}

Here's something interesting.  When we look at \emph{who} gets removed from
the top 30\% by the re-ranking:

At baseline (no correction), Causal Forest's top 30\% is approximately:
\begin{itemize}[topsep=2pt,itemsep=1pt]
  \item \textbf{8.4\%} users who genuinely benefited from the ad ($\tau = +1$)
  \item \textbf{83.5\%} users with zero effect ($\tau = 0$)---sure-things and lost causes
  \item \textbf{8.1\%} users who were hurt by the ad ($\tau = -1$)---sleeping dogs
\end{itemize}

After applying \mstpu, the sleeping dog fraction drops from 8.1\% to 3.1\%.
That's genuinely good.  But persuadables also drop from 8.4\% to 4.0\%.
The empty slots fill back up with zero-effect users.  So the mean lift of
targeted users actually improves (iROAS goes up), but wasted spend goes up
too---because the definition of ``wasted'' includes zero-effect users, and
there's now a higher fraction of them.

Figure~\ref{fig:sleeping_dog} shows this pattern across all models and datasets.

\begin{figure}[t]
  \centering
  \figinclude{figures/fig5_sleeping_dog.pdf}
  \caption{Who's in the top 30\%? Green bars are persuadables ($\tau > 0$),
    red bars are sleeping dogs ($\tau < 0$).  \mstpu\ cuts sleeping dogs but also
    cuts persuadables by a similar amount, which is why wasted spend doesn't improve.}
  \label{fig:sleeping_dog}
\end{figure}

\subsection{Why nothing can fix this: the signal just isn't there}

Figure~\ref{fig:iroas_alpha} shows iROAS over different levels of confounding.
\mstpu\ (solid lines) tracks almost exactly with the base model (dashed lines).
There's no separation.

The underlying reason: at 5\% conversion rates, the predicted effect $\hat{\tau}$
has a standard deviation of about 0.03 across users.  But distinguishing a
``definitely buys because of ad'' user from a ``probably doesn't'' user requires
more separation than that.  Since \mstpu\ can only reshuffle based on
$\hat{\tau}$ and $\hat{g}$---information we already have---it can't create
ranking signal that wasn't there to begin with.

\begin{figure}[t]
  \centering
  \figinclude{figures/fig4_iroas_by_alpha.pdf}
  \caption{iROAS efficiency at different confounding levels.
    Solid lines are \mstpu, dashed are the base model.  There's almost no
    separation---re-ranking doesn't change performance as confounding varies.}
  \label{fig:iroas_alpha}
\end{figure}

% ============================================================
\section{What This Means in Practice}
\label{sec:discussion}

\paragraph{Don't use additive STPU.}
It always makes performance worse.  The math is clear: the penalty you
subtract is on a different scale than the scores you're adjusting, so it
takes over the ranking and points you at the wrong users.

\paragraph{\mstpu\ is harmless but not helpful.}
It doesn't hurt (except in a random-noise sense), and it requires no tuning.
If you want to apply it as a default, there's no real downside.  But don't
expect it to fix wasted spend.

\paragraph{The wasted spend metric can be misleading.}
``Wasted spend'' counts both sure-things and sleeping dogs as equally bad.
But sleeping dogs are worse than sure-things---they actively decrease
conversion rates.  \mstpu\ actually does a decent job removing sleeping dogs
from the targeting set.  If you track iROAS efficiency instead of raw wasted
spend, \mstpu\ looks more useful because iROAS rewards the removal of
negative-effect users.

Our recommendation: \textbf{report iROAS efficiency as your primary metric}.
Wasted spend alone doesn't tell you whether you're wasting budget on harmless
sure-things or actively harmful sleeping dogs.

\paragraph{How to actually reduce waste.}
The real bottleneck isn't the ranking algorithm---it's data quality.
With binary (yes/no) conversions at 5\%, individual users are just too noisy
to classify reliably.  The fixes that would actually work:
\begin{enumerate}[topsep=2pt,itemsep=1pt]
  \item \textbf{Continuous outcomes.}  Model revenue per user, not just
    whether they converted.  This is far less noisy and gives CATE models
    much more to work with.
  \item \textbf{Multi-touch data.}  If you observe the same user across
    multiple exposures, you get repeated measurements that reduce uncertainty.
  \item \textbf{Longer attribution windows.}  A 24-hour conversion window
    is noisy; a 30-day window has more signal about true incrementality.
\end{enumerate}

\paragraph{One real limitation to note.}
Our experiments used synthetic data where we know the true effects.
In practice, $\hat{g}$ is estimated from whoever ended up in the control
group, which may be a biased sample if treatment wasn't random.  Under
heavy confounding, $\hat{g}$ may absorb noise from confounders and
misidentify some persuadables as sure-things.

% ============================================================
\section{Summary}
\label{sec:conclusion}

We tested a natural-seeming fix for wasted ad spend: penalize users who
already have high organic conversion rates.  Here's what we found:

\begin{itemize}[topsep=2pt,itemsep=2pt]
  \item \textbf{Additive STPU always makes things worse.}
    The penalty is on a 2--5$\times$ larger scale than the lift estimates,
    so it takes over the ranking entirely.  Users get sorted by lowest
    organic rate, which is the wrong objective.

  \item \textbf{Multiplicative STPU doesn't help either.}
    It fixes the scale problem but shows no consistent improvement---48\%
    win rate across 567 experiments (chance level).

  \item \textbf{The gap between good models and oracle can't be closed
    by re-ranking.}
    At a 5\% conversion rate, only 6.5\% of users have realized positive
    treatment effects.  Even with perfect knowledge, 78\% of your 30\%
    targeting budget is ``wasted'' by the wasted-spend definition.  The
    path from 86--92\% down to 78\% requires better estimation, not
    better re-ranking.
\end{itemize}

The takeaway for practitioners: measure iROAS efficiency (not just wasted
spend), don't apply additive STPU, and if you want to genuinely reduce waste,
invest in richer outcome data rather than better sorting algorithms.

Code and data are available at
\href{https://github.com/natepuppy/ads-uplift-benchmark}{github.com/natepuppy/ads-uplift-benchmark}.

% ============================================================
\bibliographystyle{plainnat}
\bibliography{refs}

% ============================================================
\appendix
\section{The Scale Mismatch: A Formal Statement}
\label{app:scale}

For readers who want the precise version of the scale argument:

\begin{proposition}
Let $\sigma_\tau = \mathrm{std}(\tauhat)$ and $\sigma_g = \mathrm{std}(\ghat)$,
and write $r = \sigma_g / \sigma_\tau > 1$.
The Pearson correlation between the additive STPU score
$s_\beta = \tauhat - \beta\,\ghat$ and the base score $\tauhat$ satisfies:
\begin{equation}
  \mathrm{corr}(s_\beta,\, \tauhat) \;\leq\;
    \frac{1}{\sqrt{1 + (\beta r)^2 (1 - \rho_{t,g}^2)}}
\end{equation}
where $\rho_{t,g} = \mathrm{corr}(\tauhat, \ghat)$.
\end{proposition}

In plain English: as you increase $\beta$, the STPU score ranking diverges
from the original lift-based ranking.  At $r \approx 2.5$ (what we observe
empirically) and $\beta = 0.5$, the correlation drops to 0.68---meaning about
a third of the ranking information has been replaced by ``who has the lowest
organic rate.''

The proof follows from expanding $\mathrm{Var}(s_\beta)$ and
$\mathrm{Cov}(s_\beta, \tauhat)$ using standard covariance rules.

\section{Full Beta Sweep Results}
\label{app:beta}

Figure~\ref{fig:beta_sensitivity} shows wasted spend at each tested $\beta$
value, for all DGP/model combinations.  Every single panel shows wasted spend
flat or increasing as $\beta$ grows.  There is no setting where increasing the
additive penalty helps.

\begin{figure}[H]
  \centering
  \figinclude{figures/fig3_beta_sensitivity.pdf}
  \caption{Wasted spend vs.\ $\beta$ across all datasets and models.
    The dashed horizontal line marks the base model ($\beta = 0$).
    Every panel shows the same pattern: increasing $\beta$ never helps.}
  \label{fig:beta_sensitivity}
\end{figure}

\end{document}
